\documentclass[12pt]{article}
\usepackage[utf8]{inputenc}
\usepackage[T1]{fontenc}
\usepackage{amsmath, amssymb}
\usepackage{geometry}
\usepackage{xcolor}
\usepackage{lipsum}
\usepackage{titlesec}
\usepackage{fancyhdr}
\usepackage[most]{tcolorbox}
\usepackage{graphicx}
\usepackage{float}
\usepackage{tikz}
\usetikzlibrary{shapes.geometric, arrows.meta}
\usetikzlibrary{arrows.meta, decorations.markings}

\geometry{a4paper, margin=2.5cm}
\pagestyle{fancy}
\fancyhf{}
\rhead{Problem Set VII}
\lhead{Rigid body mechanics}
\cfoot{\thepage}

% Fancy titles
\titleformat{\section}{\normalfont\Large\bfseries}{Exercise \thesection:}{1em}{}
\titleformat{\subsection}{\normalfont\bfseries}{Correction:}{1em}{}

% Custom box for correction with breakable option
\newtcolorbox{correctionbox}{
  colback=gray!5,
  colframe=black,
  fonttitle=\bfseries,
  title=Correction,
  before skip=10pt,
  after skip=10pt,
  breakable,
  enhanced
}

% Fix headheight warning
\setlength{\headheight}{14.49998pt}

\begin{document}

\begin{center}
\Large\textbf{Ibn Tofail University} \\[1em]
\large\textit{Rigid body mechanics} \\[2em]
\large\textit{Problem Set VII} \\[0.5em]
\end{center}

\vspace{1cm}

% ----------- EXERCISE 1 -------------
\section{}

A homogeneous solid sphere $S$, of mass $m$, radius $a$, and center $G$, rolls without slipping on a fixed horizontal plane $\Pi$, located at elevation $z = 0$, relative to a frame $R(O, X, Y, Z)$, with basis $(\vec{i}, \vec{j}, \vec{k})$, assumed to be Galilean. The $OZ$ axis is vertical upward. We denote by $H$ the intersection point of plane $\Pi$ with the $OZ$ axis. We denote by $I$ the geometric contact point of $S$ with $\Pi$ and $\vec{\Omega}$ the instantaneous rotation vector of $S$ with respect to $R$. The torsor of contact actions at point $I$ is reduced to the reaction $\vec{R}$.

\begin{enumerate}
    \item Write the rolling without slipping relation of $S$ on plane $\Pi$, then deduce the velocity and acceleration vectors of the center of mass $G$, with respect to $R$, as a function of $\vec{\Omega}$ and $\dot{\vec{\Omega}}$.
    \item Determine the time-dependent vector equation of motion of $G$, $\overrightarrow{HG}(t)$, by applying the general theorems. What is the nature of the motion of $G$? We denote by $\vec{v}_0$ the initial velocity vector of $G$ and by $\vec{r}_0$ the initial position vector of $G$.
\end{enumerate}

In what follows, we assume that plane $\Pi$ is rotating about the $OZ$ axis and we denote by $\vec{\omega}$ the instantaneous rotation vector of the plane with respect to $R$.

\begin{enumerate}
    \setcounter{enumi}{2}
    \item Determine the velocity of center of mass $G$, $\vec{v}_G$, using the rolling without slipping relation.
    \item Find, using the general theorems, the acceleration of center of mass $G$, $\vec{\gamma}_G$, as a function of $\vec{\omega}$ and $\vec{\Omega}$ and their time derivatives. Deduce a first vector integral of motion.
\end{enumerate}

Let $r$ and $\varphi$ be the polar coordinates of $G$, defined in the plane $XOY$ by $\overrightarrow{HG} = r\vec{u}_r$. Initially, we have the following values $r_0 = r(0)$, $\dot{r}_0 = \dot{r}(0)$, and $\dot{\varphi}_0 = \dot{\varphi}(0)$. And we denote by $\mu$ the sliding friction coefficient of $S$ on plane $\Pi$.

\begin{enumerate}
    \setcounter{enumi}{4}
    \item Determine the time-dependent equations of motion $r(t)$ and $\varphi(t)$, and deduce $\overrightarrow{HG}(t)$ as a function of time.
    \item Determine $\vec{R}$ and $\vec{M}_I$ and indicate the condition that $\mu$ must satisfy for rolling without slipping to persist.
\end{enumerate}

\begin{correctionbox}
\textbf{1. Rolling without slipping relation and velocity/acceleration of G:}

The rolling without slipping condition means the velocity of the contact point $I$ relative to the plane is zero:
$$\vec{V}_I(S/R) = \vec{0}$$

Using the velocity distribution formula:
$$\vec{V}_I(S/R) = \vec{V}_G(S/R) + \vec{\Omega} \times \vec{GI}$$

Since $\vec{GI} = -a\vec{k}$ (assuming the sphere is above the plane), we have:
$$\vec{V}_G(S/R) + \vec{\Omega} \times (-a\vec{k}) = \vec{0}$$
$$\vec{V}_G(S/R) = a\vec{\Omega} \times \vec{k}$$

This is the rolling without slipping relation.

Now, differentiating with respect to time to get acceleration:
$$\vec{\gamma}_G(S/R) = \frac{d\vec{V}_G}{dt} = a\frac{d\vec{\Omega}}{dt} \times \vec{k} + a\vec{\Omega} \times \frac{d\vec{k}}{dt}$$

Since $\vec{k}$ is fixed in the Galilean frame $R$, $\frac{d\vec{k}}{dt} = \vec{0}$, so:
$$\vec{\gamma}_G(S/R) = a\dot{\vec{\Omega}} \times \vec{k}$$

\textbf{2. Equation of motion of G and nature of motion:}

Apply the momentum theorem:
$$m\vec{\gamma}_G = m\vec{g} + \vec{R}$$

Where $\vec{R}$ is the reaction at the contact point.

The angular momentum theorem about point $G$:
$$\frac{d\vec{\sigma}_G}{dt} = \vec{M}_G(\vec{R})$$

For a sphere: $\vec{\sigma}_G = I_G\vec{\Omega} = \frac{2}{5}ma^2\vec{\Omega}$

And $\vec{M}_G(\vec{R}) = \vec{GI} \times \vec{R} = (-a\vec{k}) \times \vec{R}$

So:
$$\frac{2}{5}ma^2\dot{\vec{\Omega}} = -a\vec{k} \times \vec{R}$$
$$\dot{\vec{\Omega}} = -\frac{5}{2ma}\vec{k} \times \vec{R}$$

Now substitute into the acceleration expression:
$$\vec{\gamma}_G = a\left(-\frac{5}{2ma}\vec{k} \times \vec{R}\right) \times \vec{k} = -\frac{5}{2m}(\vec{k} \times \vec{R}) \times \vec{k}$$

Using the vector triple product identity: $(\vec{A} \times \vec{B}) \times \vec{C} = (\vec{A} \cdot \vec{C})\vec{B} - (\vec{B} \cdot \vec{C})\vec{A}$

So: $(\vec{k} \times \vec{R}) \times \vec{k} = (\vec{k} \cdot \vec{k})\vec{R} - (\vec{R} \cdot \vec{k})\vec{k} = \vec{R} - R_z\vec{k}$

Therefore:
$$\vec{\gamma}_G = -\frac{5}{2m}(\vec{R} - R_z\vec{k})$$

From the momentum theorem: $m\vec{\gamma}_G = m\vec{g} + \vec{R}$

Substitute:
$$m\left[-\frac{5}{2m}(\vec{R} - R_z\vec{k})\right] = m\vec{g} + \vec{R}$$
$$-\frac{5}{2}(\vec{R} - R_z\vec{k}) = m\vec{g} + \vec{R}$$
$$-\frac{5}{2}\vec{R} + \frac{5}{2}R_z\vec{k} = m\vec{g} + \vec{R}$$
$$-\frac{7}{2}\vec{R} + \frac{5}{2}R_z\vec{k} = m\vec{g}$$

Separate into horizontal and vertical components:

Vertical: $-\frac{7}{2}R_z + \frac{5}{2}R_z = -mg \Rightarrow -R_z = -mg \Rightarrow R_z = mg$

Horizontal: $-\frac{7}{2}\vec{R}_h = \vec{0} \Rightarrow \vec{R}_h = \vec{0}$

So the reaction is purely vertical: $\vec{R} = mg\vec{k}$

Then from the momentum theorem: $m\vec{\gamma}_G = m\vec{g} + mg\vec{k} = \vec{0}$

So $\vec{\gamma}_G = \vec{0}$, meaning the center of mass moves with constant velocity.

The equation of motion is: $\overrightarrow{HG}(t) = \overrightarrow{HG}(0) + \vec{v}_0 t$

The motion of $G$ is uniform rectilinear motion.

\textbf{3. Velocity of G with rotating plane:}

When the plane rotates with angular velocity $\vec{\omega} = \omega\vec{k}$, the rolling without slipping condition becomes that the velocity of the contact point relative to the rotating plane is zero.

Let $R'$ be the frame rotating with the plane. The condition is:
$$\vec{V}_I(S/R') = \vec{0}$$

Using the composition of velocities:
$$\vec{V}_I(S/R) = \vec{V}_I(S/R') + \vec{V}_I(R'/R)$$

But $\vec{V}_I(R'/R) = \vec{\omega} \times \overrightarrow{HI}$

And $\overrightarrow{HI} = \overrightarrow{HG} + \overrightarrow{GI} = \vec{r} - a\vec{k}$, where $\vec{r} = \overrightarrow{HG}$

So: $\vec{V}_I(R'/R) = \vec{\omega} \times (\vec{r} - a\vec{k}) = \omega\vec{k} \times \vec{r}$ (since $\vec{k} \times \vec{k} = \vec{0}$)

The rolling condition $\vec{V}_I(S/R') = \vec{0}$ gives:
$$\vec{V}_I(S/R) = \omega\vec{k} \times \vec{r}$$

But also: $\vec{V}_I(S/R) = \vec{V}_G(S/R) + \vec{\Omega} \times \vec{GI} = \vec{V}_G + \vec{\Omega} \times (-a\vec{k})$

So: $\vec{V}_G + \vec{\Omega} \times (-a\vec{k}) = \omega\vec{k} \times \vec{r}$

Therefore: $\vec{V}_G = a\vec{\Omega} \times \vec{k} + \omega\vec{k} \times \vec{r}$

\textbf{4. Acceleration of G and first integral:}

Differentiate the velocity expression:
$$\vec{\gamma}_G = \frac{d\vec{V}_G}{dt} = a\dot{\vec{\Omega}} \times \vec{k} + a\vec{\Omega} \times \frac{d\vec{k}}{dt} + \dot{\omega}\vec{k} \times \vec{r} + \omega\vec{k} \times \vec{V}_G$$

Since $\vec{k}$ is fixed and assuming constant $\omega$:
$$\vec{\gamma}_G = a\dot{\vec{\Omega}} \times \vec{k} + \omega\vec{k} \times \vec{V}_G$$

Now apply the momentum theorem:
$$m\vec{\gamma}_G = m\vec{g} + \vec{R}$$

And the angular momentum theorem about $G$:
$$\frac{d\vec{\sigma}_G}{dt} = \vec{M}_G(\vec{R}) = \vec{GI} \times \vec{R} = (-a\vec{k}) \times \vec{R}$$

With $\vec{\sigma}_G = I_G\vec{\Omega} = \frac{2}{5}ma^2\vec{\Omega}$, so:
$$\frac{2}{5}ma^2\dot{\vec{\Omega}} = -a\vec{k} \times \vec{R}$$
$$\dot{\vec{\Omega}} = -\frac{5}{2ma}\vec{k} \times \vec{R}$$

Substitute into the acceleration:
$$\vec{\gamma}_G = a\left(-\frac{5}{2ma}\vec{k} \times \vec{R}\right) \times \vec{k} + \omega\vec{k} \times \vec{V}_G$$
$$= -\frac{5}{2m}(\vec{k} \times \vec{R}) \times \vec{k} + \omega\vec{k} \times \vec{V}_G$$

Using the vector identity again: $(\vec{k} \times \vec{R}) \times \vec{k} = \vec{R} - R_z\vec{k}$

So:
$$\vec{\gamma}_G = -\frac{5}{2m}(\vec{R} - R_z\vec{k}) + \omega\vec{k} \times \vec{V}_G$$

From momentum theorem: $m\vec{\gamma}_G = m\vec{g} + \vec{R}$

So:
$$m\left[-\frac{5}{2m}(\vec{R} - R_z\vec{k}) + \omega\vec{k} \times \vec{V}_G\right] = m\vec{g} + \vec{R}$$
$$-\frac{5}{2}(\vec{R} - R_z\vec{k}) + m\omega\vec{k} \times \vec{V}_G = m\vec{g} + \vec{R}$$
$$-\frac{7}{2}\vec{R} + \frac{5}{2}R_z\vec{k} + m\omega\vec{k} \times \vec{V}_G = m\vec{g}$$

Separate components:

Vertical: $-\frac{7}{2}R_z + \frac{5}{2}R_z = -mg \Rightarrow -R_z = -mg \Rightarrow R_z = mg$

Horizontal: $-\frac{7}{2}\vec{R}_h + m\omega\vec{k} \times \vec{V}_G = \vec{0}$

So: $\vec{R}_h = \frac{2}{7}m\omega\vec{k} \times \vec{V}_G$

Now the momentum theorem gives:
$$m\vec{\gamma}_G = m\vec{g} + mg\vec{k} + \frac{2}{7}m\omega\vec{k} \times \vec{V}_G$$
$$\vec{\gamma}_G = \frac{2}{7}\omega\vec{k} \times \vec{V}_G$$

This is the equation of motion. To find a first integral, take the dot product with $\vec{V}_G$:
$$\vec{\gamma}_G \cdot \vec{V}_G = \frac{2}{7}\omega(\vec{k} \times \vec{V}_G) \cdot \vec{V}_G = 0$$

Since $(\vec{k} \times \vec{V}_G) \cdot \vec{V}_G = 0$ (perpendicular vectors).

So $\frac{d}{dt}(\frac{1}{2}V_G^2) = 0$, meaning the speed is constant: $V_G = \text{constant}$

This is the first integral.

\textbf{5. Equations of motion in polar coordinates:}

In polar coordinates: $\vec{r} = r\vec{u}_r$, $\vec{V}_G = \dot{r}\vec{u}_r + r\dot{\varphi}\vec{u}_\varphi$

The acceleration is:
$$\vec{\gamma}_G = (\ddot{r} - r\dot{\varphi}^2)\vec{u}_r + (r\ddot{\varphi} + 2\dot{r}\dot{\varphi})\vec{u}_\varphi$$

From the equation of motion: $\vec{\gamma}_G = \frac{2}{7}\omega\vec{k} \times \vec{V}_G$

Now $\vec{k} \times \vec{V}_G = \vec{k} \times (\dot{r}\vec{u}_r + r\dot{\varphi}\vec{u}_\varphi) = \dot{r}\vec{u}_\varphi - r\dot{\varphi}\vec{u}_r$

So:
$$\vec{\gamma}_G = \frac{2}{7}\omega(\dot{r}\vec{u}_\varphi - r\dot{\varphi}\vec{u}_r) = -\frac{2}{7}\omega r\dot{\varphi}\vec{u}_r + \frac{2}{7}\omega\dot{r}\vec{u}_\varphi$$

Equate components:

Radial: $\ddot{r} - r\dot{\varphi}^2 = -\frac{2}{7}\omega r\dot{\varphi}$

Angular: $r\ddot{\varphi} + 2\dot{r}\dot{\varphi} = \frac{2}{7}\omega\dot{r}$

These are the equations of motion.

From the first integral, we know $V_G^2 = \dot{r}^2 + r^2\dot{\varphi}^2 = \text{constant}$

The position vector is: $\overrightarrow{HG}(t) = r(t)\vec{u}_r(\varphi(t))$

\textbf{6. Reaction $\vec{R}$ and moment $\vec{M}_I$, condition on $\mu$:}

We already found:
$$\vec{R} = \vec{R}_h + R_z\vec{k} = \frac{2}{7}m\omega\vec{k} \times \vec{V}_G + mg\vec{k}$$

For rolling without slipping to persist, the friction force must satisfy the Coulomb condition:
$$\|\vec{R}_h\| \leq \mu R_z = \mu mg$$

So:
$$\left\|\frac{2}{7}m\omega\vec{k} \times \vec{V}_G\right\| \leq \mu mg$$
$$\frac{2}{7}\omega V_G \leq \mu g$$
$$\mu \geq \frac{2\omega V_G}{7g}$$

This is the condition on $\mu$.

The moment at point $I$ is zero for a perfect rolling without slipping contact (no friction moment for a sphere on a plane).

Therefore, the torsor of contact actions at $I$ is:
$$\tau(I) = [\vec{R}, \vec{0}]$$
with $\vec{R} = \frac{2}{7}m\omega\vec{k} \times \vec{V}_G + mg\vec{k}$
\end{correctionbox}

\end{document}